\section*{Capítulo \ref{ch:g2:caring-for-nature} --- Cuidar da natureza}

\begin{solution}{exo:g2:caring-for-nature:1}
Aquilo que está em volta de um \omterm{def:g1:living-or-not:living}{ser vivo} e do que ele depende. Por
exemplo: o ar que respiramos, a água que bebemos, o alimento que
comemos --- tudo vindo da \omterm{def:g2:caring-for-nature:environment}{natureza}.
\end{solution}

\begin{solution}{exo:g2:caring-for-nature:2}
Por exemplo: os besouros debaixo da casca, os tatuzinhos por baixo
do tronco, o musgo em cima, o ouriço passando o inverno no oco.
\end{solution}

\begin{solution}{exo:g2:caring-for-nature:3}
Por exemplo: levar todo o lixo para casa; olhar em vez de coletar, e
nunca colher um tufo inteiro; recolocar cada pedra ou tronco
virado; ficar nas trilhas perto de lugares frágeis; observar os
\omterm{def:g1:animals-around-us:animal}{animais} de longe, em silêncio, e não mexer nos filhotes.
\end{solution}

\begin{solution}{exo:g2:caring-for-nature:4}
Porque a \omterm{def:g2:caring-for-nature:environment}{natureza} não tem caminhão de lixo: o que deixamos, fica ---
uma tampinha de plástico por muitos e muitos anos, tempo suficiente
para prender ou envenenar \omterm{def:g1:animals-around-us:animal}{animais} e ainda estar ali quando as
crianças de hoje forem adultas.
\end{solution}

\begin{solution}{exo:g2:caring-for-nature:5}
Colhido até a última \omterm{def:g1:plants-around-us:parts}{flor}, o tufo não consegue fazer \omterm{def:g2:from-seed-to-plant:seed}{sementes} ---
então não haverá margaridas ali no ano que vem --- e as abelhas
perdem as \omterm{def:g1:plants-around-us:parts}{flores} de que se alimentam.
\end{solution}

\begin{solution}{exo:g2:caring-for-nature:6}
Eles apodrecem devagar até virar uma terra escura e fofa. As
minhocas, os tatuzinhos e \omterm{def:g1:living-or-not:living}{seres vivos} pequenos demais para enxergar
fazem o trabalho; no fim a gente ganha o composto --- uma boa terra
que alimenta as \omterm{def:g1:plants-around-us:plant}{plantas} do ano seguinte.
\end{solution}

\begin{solution}{exo:g2:caring-for-nature:7}
Por exemplo: um pratinho de água numa onda de calor, um cantinho
selvagem sem corte no quintal, \omterm{def:g2:from-seed-to-plant:seed}{sementes} para as aves no frio, cavar
uma lagoinha.
\end{solution}

\begin{solution}{exo:g2:caring-for-nature:8}
Resposta aberta. Embalagens e garrafas de plástico costumam ganhar.
Cada tipo deveria ter ido para casa com o dono --- e depois para a
lixeira certa de separação.
\end{solution}

\begin{solution}{exo:g2:caring-for-nature:9}
Deixe os \omterm{ex:g2:animals-grow-up:frog}{girinos} na lagoa: ela tem tudo de que eles precisam --- a
água da lagoa, o alimento, espaço para criar \omterm{def:g1:my-body:parts}{pernas} --- e um pote
não tem. Se a Ema quiser ajudar, ela pode observar os \omterm{ex:g2:animals-grow-up:frog}{girinos} ali
durante a primavera, manter a margem sem pisoteio e talvez ajudar a
fazer um novo \omterm{def:g2:where-animals-live:habitat}{hábitat} de lagoa.
\end{solution}

\begin{solution}{exo:g2:caring-for-nature:10}
Um milhão de ``nadas'' no sentido ruim é um campo enterrado em
tampinhas --- os danos pequenos se multiplicam. No sentido bom, um
milhão de cuidados pequenos --- lixo levado para casa, torneiras
fechadas, lixo separado --- somam com a mesma certeza: a \omterm{def:g2:caring-for-nature:environment}{natureza} é
cuidada principalmente em escolhas exatamente desse tamanho.
\end{solution}
